% CI/CD, Azure DevOps & Terraform section

% 1 – What is CI/CD?
\QuestionSlide[\CategoryBadge[CICDColor!20]{CI/CD \& IaC}]{* What is CI/CD?}

\begin{frame}[fragile]
  \frametitle{%
    \begin{tikzpicture}[remember picture, overlay]
      \node[anchor=north east, xshift=-0.4cm, yshift=-0.4cm, text=black] at (current page.north east) {
        \CategoryBadge[CICDColor!20]{CI/CD \& IaC}
      };
    \end{tikzpicture}
    Answer \theqcounter: What is CI/CD?%
  }

  {\footnotesize
  \textbf{CI (Continuous Integration)} automates build and test on every commit, catching integration bugs early.
  \textbf{CD (Continuous Delivery/Deployment)} automatically delivers validated builds to staging/production. Together they shorten feedback loops, reduce manual errors, and enable multiple daily releases with high confidence.
  }
\end{frame}

% 2 – Azure DevOps Pipeline
\QuestionSlide[\CategoryBadge[CICDColor!20]{CI/CD \& IaC}]{** What is an Azure DevOps YAML Pipeline?}

\begin{frame}[fragile]
  \frametitle{%
    \begin{tikzpicture}[remember picture, overlay]
      \node[anchor=north east, xshift=-0.4cm, yshift=-0.4cm, text=black] at (current page.north east) {
        \CategoryBadge[CICDColor!20]{CI/CD \& IaC}
      };
    \end{tikzpicture}
    Answer \theqcounter: What is an Azure DevOps YAML Pipeline?%
  }

  {\footnotesize
  An \textbf{Azure Pipelines} YAML file (\texttt{azure-pipelines.yml}) defines the entire CI/CD flow as code: \textbf{stages} $\to$ \textbf{jobs} $\to$ \textbf{steps}. It runs on Microsoft-hosted or self-hosted agents. Pipeline templates enable reuse across repos and projects.
  }

  \begin{minted}[fontsize=\scriptsize]{yaml}
trigger:
  branches:
    include: [main]

stages:
- stage: Build
  jobs:
  - job: BuildAndTest
    pool:
      vmImage: ubuntu-latest
    steps:
    - task: DotNetCoreCLI@2
      inputs:
        command: test
        projects: '**/*.Tests.csproj'
    - task: DotNetCoreCLI@2
      inputs:
        command: publish
        arguments: '-c Release -o $(Build.ArtifactStagingDirectory)'
  \end{minted}
\end{frame}

% 3 – GitHub Actions
\QuestionSlide[\CategoryBadge[CICDColor!20]{CI/CD \& IaC}]{** What is a GitHub Actions workflow?}

\begin{frame}[fragile]
  \frametitle{%
    \begin{tikzpicture}[remember picture, overlay]
      \node[anchor=north east, xshift=-0.4cm, yshift=-0.4cm, text=black] at (current page.north east) {
        \CategoryBadge[CICDColor!20]{CI/CD \& IaC}
      };
    \end{tikzpicture}
    Answer \theqcounter: What is a GitHub Actions workflow?%
  }

  {\footnotesize
  GitHub Actions workflows (\texttt{.github/workflows/*.yml}) are event-driven pipelines. \textbf{Triggers}: \texttt{push}, \texttt{pull\_request}, \texttt{schedule}, or \texttt{workflow\_dispatch}. \textbf{Jobs} run in parallel by default; \texttt{needs} declares ordering. Reusable workflows reduce duplication across repos.
  }

  \begin{minted}[fontsize=\scriptsize]{yaml}
on:
  push:
    branches: [main]

jobs:
  build-and-push:
    runs-on: ubuntu-latest
    steps:
    - uses: actions/checkout@v4
    - uses: actions/setup-dotnet@v4
      with:
        dotnet-version: '9.x'
    - run: dotnet test
    - run: docker build -t ghcr.io/${{ github.repository }}/api:${{ github.sha }} .
    - run: docker push ghcr.io/${{ github.repository }}/api:${{ github.sha }}
  \end{minted}
\end{frame}

% 4 – Terraform
\QuestionSlide[\CategoryBadge[CICDColor!20]{CI/CD \& IaC}]{** What is Terraform?}

\begin{frame}[fragile]
  \frametitle{%
    \begin{tikzpicture}[remember picture, overlay]
      \node[anchor=north east, xshift=-0.4cm, yshift=-0.4cm, text=black] at (current page.north east) {
        \CategoryBadge[CICDColor!20]{CI/CD \& IaC}
      };
    \end{tikzpicture}
    Answer \theqcounter: What is Terraform?%
  }

  {\footnotesize
  \textbf{Terraform} is HashiCorp's open-source \textbf{Infrastructure as Code} tool. You write declarative HCL (HashiCorp Configuration Language) describing the desired cloud state; Terraform diffs it against the current state and generates a \textbf{plan} showing exactly what will change before applying.
  }

  \begin{minted}[fontsize=\scriptsize]{hcl}
resource "azurerm_app_service" "api" {
  name                = "my-api-prod"
  location            = azurerm_resource_group.rg.location
  resource_group_name = azurerm_resource_group.rg.name
  app_service_plan_id = azurerm_app_service_plan.plan.id

  app_settings = {
    ASPNETCORE_ENVIRONMENT = "Production"
  }
}
  \end{minted}
\end{frame}

% 5 – Terraform State
\QuestionSlide[\CategoryBadge[CICDColor!20]{CI/CD \& IaC}]{** What is Terraform state and why does it need remote storage?}

\begin{frame}[fragile]
  \frametitle{%
    \begin{tikzpicture}[remember picture, overlay]
      \node[anchor=north east, xshift=-0.4cm, yshift=-0.4cm, text=black] at (current page.north east) {
        \CategoryBadge[CICDColor!20]{CI/CD \& IaC}
      };
    \end{tikzpicture}
    Answer \theqcounter: What is Terraform state?%
  }

  {\footnotesize
  \textbf{State} (\texttt{terraform.tfstate}) maps HCL resources to real cloud IDs. Without it, Terraform cannot know what already exists. Store state \textbf{remotely} (Azure Blob, Terraform Cloud, S3) so the whole team shares the same view and concurrent runs are serialised by \textbf{state locking}.
  }

  \begin{minted}[fontsize=\scriptsize]{hcl}
terraform {
  backend "azurerm" {
    resource_group_name  = "tf-state-rg"
    storage_account_name = "tfstatecontoso"
    container_name       = "tfstate"
    key                  = "prod.terraform.tfstate"
  }
}
  \end{minted}
\end{frame}

% 6 – Blue-Green vs Canary
\QuestionSlide[\CategoryBadge[CICDColor!20]{CI/CD \& IaC}]{** What is the difference between Blue-Green and Canary deployments?}

\begin{frame}[fragile]
  \frametitle{%
    \begin{tikzpicture}[remember picture, overlay]
      \node[anchor=north east, xshift=-0.4cm, yshift=-0.4cm, text=black] at (current page.north east) {
        \CategoryBadge[CICDColor!20]{CI/CD \& IaC}
      };
    \end{tikzpicture}
    Answer \theqcounter: Blue-Green vs Canary deployments?%
  }

  {\footnotesize
  \textbf{Blue-Green}: two identical environments — Blue (live), Green (new version). Traffic swaps instantly via a load balancer rule; rollback means switching back. Zero downtime but requires double the infrastructure.
  \vspace{0.5em}

  \textbf{Canary}: the new version receives a small percentage of traffic (e.g.\ 5\%); metrics are monitored; if healthy, traffic gradually shifts to 100\%. Slower but detects production issues with minimal user impact.
  }
\end{frame}

% 7 – Environment Protection Rules
\QuestionSlide[\CategoryBadge[CICDColor!20]{CI/CD \& IaC}]{** What are environment protection rules in GitHub Actions / Azure DevOps?}

\begin{frame}[fragile]
  \frametitle{%
    \begin{tikzpicture}[remember picture, overlay]
      \node[anchor=north east, xshift=-0.4cm, yshift=-0.4cm, text=black] at (current page.north east) {
        \CategoryBadge[CICDColor!20]{CI/CD \& IaC}
      };
    \end{tikzpicture}
    Answer \theqcounter: What are environment protection rules?%
  }

  {\footnotesize
  \textbf{Protection rules} gate deployments to critical environments (staging, production): required approvals from reviewers, branch filters (only \texttt{main} can deploy to prod), wait timers, and secret scoping. In GitHub Actions, environments are declared on a job with \texttt{environment: production}; Azure DevOps uses \textbf{Environments} with approval checks.
  }

  \begin{minted}[fontsize=\scriptsize]{yaml}
jobs:
  deploy-prod:
    environment:
      name: production
      url: https://api.contoso.com
    runs-on: ubuntu-latest
    steps:
    - run: echo "Deploying to production after approval"
  \end{minted}
\end{frame}

% 8 – Azure DevOps Artifacts
\QuestionSlide[\CategoryBadge[CICDColor!20]{CI/CD \& IaC}]{* What are Azure DevOps Artifacts?}

\begin{frame}[fragile]
  \frametitle{%
    \begin{tikzpicture}[remember picture, overlay]
      \node[anchor=north east, xshift=-0.4cm, yshift=-0.4cm, text=black] at (current page.north east) {
        \CategoryBadge[CICDColor!20]{CI/CD \& IaC}
      };
    \end{tikzpicture}
    Answer \theqcounter: What are Azure DevOps Artifacts?%
  }

  {\footnotesize
  \textbf{Azure DevOps Artifacts} is a private package feed supporting NuGet, npm, Maven, PyPI, and Universal Packages. Teams publish internal libraries here so they can be consumed like public packages. \textbf{Upstream sources} proxy public registries (nuget.org) through the same feed, enabling policy control and caching.
  }

  \begin{minted}[fontsize=\scriptsize]{xml}
<!-- nuget.config pointing to internal feed -->
<configuration>
  <packageSources>
    <add key="internal" value=
  "https://pkgs.dev.azure.com/contoso/myproject/_packaging/main/nuget/v3/index.json"/>
  </packageSources>
</configuration>
  \end{minted}
\end{frame}
